\subsection{A regime map of asset responses}

\begin{qfdefinition}[A regime map of asset responses: growth, recession, inflation, deflation, liquidity and geopolitical shocks across rates, credit, equities and sectors, currencies and commodities, as hypotheses to test]\label{def:macro-regime-asset-map}
\qffield{Statement} A \emph{regime map} lists the direction in which each asset class is expected to move under a type of macro shock. Gliner's version is qualitative: a risk-on/risk-off grouping of assets, correlations that rise in stress (Figure~4.5, around Lehman), and the need to know which relationships hold in the current environment \citep[chap.~4, sec.\ ``Assessing the Relationships between Assets'', PDF pp.~98--102]{s7-gliner}; sector rotation, defensive sectors such as consumer staples in recession, and the need to know ``whether or not the economy is in contraction, recovery, expansion, or slowing growth'' (chap.~8, sec.\ ``Sector Rotation'', PDF pp.~216--217, Figure~8.3). The map, as hypotheses to test:
\begin{itemize}
\item \emph{Growth shock}: yields up, curve steeper, credit tighter, equities, cyclicals and industrial commodities up, growth-linked currencies (AUD, CAD, EM) up.
\item \emph{Recession}: yields down, curves bull-steepen or flatten, credit wider, equities down, USD up, commodities down except gold; defensives (staples, utilities, health care) lead.
\item \emph{Inflation shock}: nominal yields up, curve bear-steepens, TIPS beat nominals, commodities (energy, agriculture) up; equities depend on real rates; low-credibility central banks' currencies punished.
\item \emph{Deflation or disinflation}: nominal yields and breakevens down, duration rallies.
\item \emph{Liquidity shock or financial crisis}: USD up sharply, Treasuries rally, credit spreads gap wider, correlations jump, volatility spikes, the cross-currency basis widens, funding markets seize.
\item \emph{USD weakness}: commodity prices up. \emph{Geopolitical crisis}: energy and gold up.
\end{itemize}
\emph{Sector rotation against Gliner's Figure~8.3}. The figure lists four phases: contraction (health care, telecom, consumer staples, utilities), recovery (financials, materials, consumer discretionary, technology), expansion (industrials, energy, materials, technology) and slowing growth (financials, telecom, consumer staples, utilities). An earlier sector-rotation table in these notes matched it only in part: its recession list (staples, utilities, health care) omitted telecom; its recovery list added real estate, which the figure does not show, and omitted materials and technology; its expansion list omitted energy; and its late-cycle ``Energy, Materials (inflation hedges)'' contradicts the figure's slowing-growth phase, which is defensive. Both are schematic: Gliner's figure gives no data, period or source.

\qffield{Notation} A \emph{bull steepener} is a steeper curve with falling yields, a \emph{bear steepener} with rising yields (\hyperref[meth:curve-trades-carry-rolldown]{curve trades}).

\qffield{Conditions} None of these signs is a law. They depend on the monetary regime and on what kind of shock moves inflation: the stock--bond correlation, for example, has changed sign across decades. Classifying regimes needs information available at the time; a map built on ex post NBER dates or revised data is look-ahead.

\qffield{Practice} Use the map to state a thesis's expressions and its hedges, then test each sign on dated regimes with real-time data before relying on it. Reading the curve as a summary of the macro narrative, or fixed income as telling what the market expects before it happens, rests on the expectations component of yields, which also contain term premia (\hyperref[mod:forward-rates-fra]{forward rates and term-structure theories}).

\qffield{Pitfalls} No source here supports an ordering such as real rates falling first in an inflation shock, or the curve bull-steepening and then flattening in a recession; the ordering depends on the policy response. An inverse correlation between commodities and the dollar holds on average in some samples and fails in others. \emph{Own test of the recession row in real time}. The hypothesis and pass rule were fixed in the script before its first run, though not preregistered. Regimes were classified with a Sahm-type rule on each quarterly vintage of the Philadelphia Fed's Real-Time Data Set, 1965--2025: the three-month average unemployment rate at least 0.5 point above its minimum over the previous twelve months. That gives 241 vintages; the three 2026 vintages are skipped because October 2025 is missing in them. Asset responses are measured over the next three months.

\emph{Real-time test.} The rule signals in 53 quarters, from 10 episodes. None of the recession signs is supported at the pre-set $\abs{t}>2$ (Welch).
\begin{itemize}
\item The 10-year yield change is $-0.05$ point relative to other quarters ($t=-0.5$), and the Baa spread change $-0.07$ ($t=-0.8$).
\item Equity excess returns are $+1.2$ percent ($t=+0.8$). The sign is wrong, since the rule fires late and recoveries follow.
\item Defensive sectors (staples, utilities, health care) less the market return $-0.7$ ($t=-1.0$), also the wrong sign.
\item The dollar is the closest. The Fed's broad index weighted goods trade only until its 2019 revision, which added services, simplified the weights and set the new index from January 2006 \citep[secs.~II--III]{fed-dollar-index-revisions-2019}. On the goods-only broad index, a consistent definition for 1973--2019 and now the primary dollar measure, it rises $+1.1$ percent ($t=+1.98$, Newey--West $t=+1.87$, 6 episodes): the right sign, just short of the threshold. On the goods-and-services index from 2006 it moves $-0.4$ ($t=-0.3$, 3 episodes).
\end{itemize}

\emph{Ex-post description.} NBER windows are dated afterwards, so they describe what happens inside recessions; they are not a usable classification. Under the \emph{any-overlap} definition, a quarter counts if any of its three response months lies inside a recession. On that definition the signs are as the map says: yields $-0.16$, spread $+0.25$, equities $-3.1$ percent, defensives $+1.1$, dollar $+0.8$ (on a splice of the two indices). None reaches Welch $\abs{t}=2$ (spread 1.78), though the spread does on a Newey--West $t$ with four lags ($2.46$). The \emph{all-months} definition, which requires all three months inside a recession, is a sensitivity, not chosen for significance. On it the yield ($-0.41$, $t=-2.2$) and spread ($+0.45$, $t=+2.9$) signs pass Welch $\abs{t}>2$; defensives ($+2.1$, $t=1.9$; Newey--West $2.7$), equities and the dollar do not.

Only with the any-overlap windows and Welch $t$ does no ex-post sign reach $\abs{t}=2$. Signal and recession quarters cluster in 6--10 episodes, so every $t$ here overstates the independent evidence. The recession row describes recessions dated afterwards; it is not a tradable real-time rule.

\emph{A second row: inflation shock}. The protocol was written before the data were downloaded, with two source substitutions recorded before computing; it is not a preregistration. The protocol names the Philadelphia Fed's quarterly CPI vintages. The quarterly file, \texttt{cpiQvMd.xlsx}, exists \citep[CPI data-files page, dated 2 October 2026; specific notes for CPI and PCPI; general notes, sec.~II and IV]{philfed-rtdsm-cpi}, and its timing differs: a quarterly vintage holds what was available on the 15th of the quarter's middle month and excludes any release published after the 15th, while a monthly vintage holds what was available just after the BLS's monthly report, which may come a few days before or after the 15th. A quarter signals an inflation shock when twelve-month CPI inflation in that vintage's latest month exceeds its level a year earlier by at least one point.

The hypothesised responses over the next three months are 10-year yields up, the 10-year--3-month slope up (bear steepening) and WTI spot up. The bill rate and WTI are from H.15 and EIA.
On the quarterly vintages, as the protocol specified:
\begin{itemize}
\item 128 vintages, 1994Q3 to 2026Q2 (2026Q3 skipped: responses incomplete); 34 signal quarters in 13 episodes.
\item 10-year yield: $+0.07$ point against other quarters, Welch $t=0.83$ (Newey--West 0.78).
\item Slope: $-0.17$, $t=-1.56$ (Newey--West $-1.49$). The opposite sign: the curve tended to flatten after accelerations.
\item WTI spot: $-2.2$ percent, $t=-0.60$.
\end{itemize}
\emph{Timing sensitivity: monthly middle-month vintages}: 111 vintages, November 1998 to May 2026; 31 signal quarters in 14 episodes; 10-year yield $+0.06$ ($t=0.76$), slope $-0.16$ ($t=-1.44$), WTI $-1.0$ percent ($t=-0.24$). On the 111 vintage months that have responses in both files, the latest CPI month differs in 55, always later in the monthly vintage (the BLS report had come out after the 15th), and the signal label differs in 11. The verdict does not change. No sign of the inflation row is supported in real time on this rule.

\emph{A third row: deflation or disinflation}. The mirror-image rule, a fall of at least one point in twelve-month CPI inflation over a year in each quarterly vintage, flags 28 of 128 quarters in 12 episodes. The hypothesised fall in 10-year yields over the next three months is not supported: $+0.04$ point against other quarters (Welch $t=0.51$, Newey--West 0.50), the wrong sign; the slope ($+0.11$, $t=1.20$) and WTI ($+5.5$ percent, $t=1.45$) also rose, reported alongside only. Limits: three rows of the map, one real-time rule each, three-month windows, U.S. assets only; CPI vintages from 1994 only; WTI spot is a price, not an investable return; TIPS, equities and currencies untested.

\qffield{See also} \hyperref[def:global-macro-strategy]{global macro}; \hyperref[def:macroeconomic-indicators]{macroeconomic indicators}; \hyperref[def:central-bank-policy-framework]{central-bank policy}; \hyperref[def:tips-breakevens]{TIPS and breakevens}.

\qffield{Further reading} Regimes and the stock--bond correlation \citep{wide-ang-timmermann-regimes}.

\end{qfdefinition}
